\documentclass[12pt]{article}
\usepackage[utf8]{inputenc}
\usepackage[T1]{fontenc}
\usepackage{hyperref}
\title{Tausug Language Resources for K-12 Curriculum Development}
\author{Compiled by Claude Sonnet 5 (claude-sonnet-5) for Andr\'as Kornai}
\date{}
\begin{document}
\maketitle

\section{Introduction}
Tausug (Bahasa S\=ug, ISO 639-3 tsg) is a language of wider communication in the southern Philippines with about 1,100,000 speakers. It is one of the 12 languages named in DepEd Order No.\ 16, s.\ 2012 for use as a learning area and medium of instruction under MTB-MLE. The searches found dictionaries and policy context, but not the actual MTB-MLE textbooks.

\section{Primary Educational Resources}
A publisher's listing for Bangahan's dictionary says its orthography is the basis of the one Bangahan wrote for the Department of Education's MTB-MLE materials for Bahasa S\=ug: \url{https://artbooks.ph/products/english-bahasa-sug-dictionary}. The MTB-MLE textbooks themselves were not found.

\section{Linguistic and Reference Materials}
\begin{itemize}
\item Bangahan, English--Bahasa S\=ug dictionary (Vibal Group, 2015), more than 15,000 entries: \url{https://artbooks.ph/products/english-bahasa-sug-dictionary}
\item Hassan, Ashley and Ashley, \emph{Tausug--English Dictionary} (Notre Dame of Jolo College, 2nd ed., 1994; Glottolog record): \url{https://glottolog.org/resource/reference/id/158613}
\item Copet, \emph{Tausug Dictionary} (1957; Glottolog record): \url{https://glottolog.org/resource/reference/id/318675}
\item Preface to the 1975 edition of the Tausug dictionary on Webonary: \url{https://www.webonary.work/tausug/overview/preface-1975-edition/}
\item DLI Tausug Language Course: \url{https://www.fsi-language-courses.org/tausug/courses/dli-tausug-language-course}
\end{itemize}

\section{Religious and Community Resources}
No Tausug Bible or other scripture translation was confirmed in the searches. This remains to be checked with Bible Society or SIL sources.

\section{Digital and Audio Resources}
The Webonary dictionary site (above) is a digital resource. The DLI course is a language-course resource whose audio content I did not verify.

\section{Curriculum Development Context}
\begin{itemize}
\item MTB-MLE policy: the DepEd Order listing is discussed in a journal article at \url{https://journals.upd.edu.ph/index.php/archive/article/download/9334/8248} and in slides at \url{https://www.slideshare.net/slideshow/knowing-about-mtbmle/68071062}.
\end{itemize}
Materials for MTB-MLE languages had to meet ``four minima'', including an orthography acceptable to the stakeholders. Obtaining the DepEd Tausug materials is the next step. Upper-grade materials were not found.
\end{document}